\documentclass[12pt,a4paper]{article}

\usepackage[T2A]{fontenc}
\usepackage[utf8]{inputenc}
\usepackage[russian,english]{babel}
\usepackage{amsmath}
\usepackage{amssymb}
\usepackage{booktabs}
\usepackage{longtable}
\usepackage{array}
\usepackage{graphicx}
\usepackage{pdflscape}
\usepackage{geometry}
\usepackage{hyperref}
\usepackage{xcolor}

\geometry{margin=2.2cm}
\setlength{\emergencystretch}{3em}
\hypersetup{
    colorlinks=true,
    linkcolor=blue!50!black,
    urlcolor=blue!50!black
}

\newcommand{\Prpow}{P_{\mathrm r}}
\newcommand{\Psbf}{P_{\mathrm{SBF}}}
\newcommand{\Pfluc}{P_{\mathrm{fluc}}}
\newcommand{\Igal}{I_{\mathrm{gal}}}
\newcommand{\Phigc}{\phi_{\mathrm{GC}}}
\newcommand{\Phibg}{\phi_{\mathrm{bg}}}

\title{Методические заметки по GO-3055:
поправка \texorpdfstring{$P_{\mathrm r}$}{Pr} и пересборка
SBF-пайплайна после курсовой работы}
\author{Рабочая методическая заметка для проекта GO-3055}
\date{27 июля 2026 г.}

\begin{document}

\maketitle

\begin{abstract}
В заметке с нуля разбирается физический смысл поправки
\(\Prpow\), её связь с измеряемой амплитудой \(P_0\), построение
функций светимости шаровых скоплений и фоновых галактик,
переход от звёздной величины к потоку, появление квадрата
потока в интеграле дисперсии и конкретная реализация в
текущем пайплайне \texttt{sbf-2.ipynb}. Отдельная часть
фиксирует все принципиальные изменения относительно версии
курсовой работы, численно сопоставляет старые и новые
измерения для всех 14 галактик GO-3055 и перечисляет
нерешённые проблемы новой реализации.
\end{abstract}

\tableofcontents

\part{Поправка \(P_{\mathrm r}\) за неразрешённые источники}

\section{Короткое определение}

\(\Prpow\) --- это ожидаемая мощность флуктуаций от компактных
источников, которые оказались слабее предела обнаружения
каталога и поэтому не были замаскированы по отдельности.

В нашей реализации учитываются две популяции:
\begin{equation}
    \Prpow = P_{\mathrm{r,GC}} + P_{\mathrm{r,bg}},
\end{equation}
где \(P_{\mathrm{r,GC}}\) --- вклад неразрешённых шаровых
скоплений, а \(P_{\mathrm{r,bg}}\) --- вклад неразрешённых
фоновых галактик.

Измеряемая по спектру мощности амплитуда \(P_0\) содержит
как настоящий SBF, так и эти источники:
\begin{equation}
    P_0 = \Psbf + \Prpow.
\end{equation}
Поэтому искомая флуктуационная мощность определяется как
\begin{equation}
    \Pfluc = \Psbf = P_0-\Prpow.
\end{equation}

Это равенство является модельным предположением. В реальном
кадре возможна дополнительная PSF-подобная мощность от пыли,
дефектов и ошибок модели галактики:
\begin{equation}
    P_0 =
    \Psbf+\Prpow+P_{\mathrm{other}}.
\end{equation}
Предполагается, что \(P_{\mathrm{other}}\) подавлено моделью,
маской и выбором рабочей области.

\section{Откуда физически возникает SBF}

Даже в двух соседних пикселях с одинаковой средней
поверхностной яркостью находится разное случайное число
звёзд. Кроме того, свет распределён между звёздами крайне
неравномерно: несколько ярких красных гигантов могут влиять
на поток сильнее, чем множество слабых звёзд.

Изображение можно представить как
\begin{equation}
    I(x,y) =
    I_{\mathrm{model}}(x,y)
    + \delta I_{\mathrm{SBF}}(x,y)
    + \delta I_{\mathrm{sources}}(x,y)
    + \delta I_{\mathrm{noise}}(x,y).
\end{equation}
Здесь \(I_{\mathrm{model}}\) --- гладкая модель галактики,
\(\delta I_{\mathrm{SBF}}\) --- искомые звёздные флуктуации,
\(\delta I_{\mathrm{sources}}\) --- компактные источники, а
\(\delta I_{\mathrm{noise}}\) --- инструментальный и фотонный
шум.

Обычные неразрешённые звёзды галактики являются искомым
сигналом. Их функцию светимости нельзя включать в
\(\Prpow\), иначе из измерения будет вычтен сам SBF.

\section{Почему компактные источники попадают именно в \(P_0\)}

\subsection{Один точечный источник}

Пусть источник с потоком \(f_i\) расположен в точке \(x_i\).
Его изображение после прохождения через телескоп:
\begin{equation}
    I_i(x) = f_i\,\mathrm{PSF}(x-x_i).
\end{equation}
Преобразование Фурье имеет вид
\begin{equation}
    \widetilde I_i(k) =
    f_i\widetilde{\mathrm{PSF}}(k)
    \exp(-ikx_i).
\end{equation}
Спектр мощности:
\begin{equation}
    \left|\widetilde I_i(k)\right|^2 =
    f_i^2
    \left|\widetilde{\mathrm{PSF}}(k)\right|^2.
\end{equation}
Положение источника исчезает из мощности, поскольку фазовый
множитель имеет единичный модуль.

\subsection{Много случайно расположенных источников}

Для совокупности независимых источников перекрёстные фазовые
члены в среднем взаимно уничтожаются, и получается
\begin{equation}
    P_{\mathrm{sources}}(k)
    \propto
    \left|\widetilde{\mathrm{PSF}}(k)\right|^2
    \sum_i f_i^2.
\end{equation}
Настоящий SBF также создаётся неразрешёнными точечными
звёздами и имеет ту же PSF-подобную форму.

Поэтому наблюдаемый спектр мощности описывается моделью
\begin{equation}
    P(k)=P_0E(k)+P_1,
\end{equation}
где:
\begin{itemize}
    \item \(P(k)\) --- измеренный азимутально усреднённый
    спектр мощности;
    \item \(E(k)\) --- ожидаемая форма спектра PSF с учётом
    рабочего кольца и маски;
    \item \(P_0\) --- амплитуда всей PSF-подобной мощности;
    \item \(P_1\) --- шумовой член, который в классической
    модели считается постоянным по \(k\).
\end{itemize}

Из одного только фита \(P(k)\) невозможно разделить
\(\Psbf\) и \(\Prpow\): обе величины умножаются на одну и ту
же функцию \(E(k)\). Поэтому \(\Prpow\) оценивается независимо
по каталогу компактных источников.

\section{Маскирование и предел обнаружения}

Звёздная величина устроена так, что меньшее \(m\) соответствует
более яркому источнику. Пусть каталог надёжен до величины
\(m_{\mathrm{lim}}\). Тогда:
\begin{align}
    m \leq m_{\mathrm{lim}}
    &\quad\Longrightarrow\quad
    \text{источник обнаружен и замаскирован},\\
    m > m_{\mathrm{lim}}
    &\quad\Longrightarrow\quad
    \text{источник остаётся в остаточном изображении}.
\end{align}

Маска и \(\Prpow\) являются двумя дополняющими частями одной
процедуры:
\begin{center}
\begin{tabular}{ll}
\toprule
Яркость источника & Действие \\
\midrule
\(m\leq m_{\mathrm{lim}}\) & индивидуальное обнаружение и маскирование \\
\(m>m_{\mathrm{lim}}\) & статистический учёт через \(\Prpow\) \\
\bottomrule
\end{tabular}
\end{center}

В текущем коде используется один порог
\texttt{source\_mask\_m\_cut}. Он задаётся вручную или
приближённо оценивается по каталогу. Строгая оценка
полноты должна в дальнейшем выполняться с помощью
искусственных источников.

\section{Функция светимости}

\subsection{Определение}

Каталог найденных источников разбивается на интервалы по
звёздной величине. Число источников нормируется на площадь
области \(N_{\mathrm{pix}}\) и ширину интервала \(dm\):
\begin{equation}
    \phi(m)=\frac{dN}{N_{\mathrm{pix}}\,dm}.
\end{equation}
Единицы \(\phi(m)\):
\begin{equation}
    [\phi] =
    \text{источников}\,
    \mathrm{pixel}^{-1}\,
    \mathrm{mag}^{-1}.
\end{equation}

Разделение шаровых скоплений и фоновых галактик выполняется
статистически. Не требуется заранее классифицировать каждый
найденный объект.

\subsection{Шаровые скопления}

Для шаровых скоплений принимается гауссовская функция
светимости GCLF:
\begin{equation}
    \Phigc(m)=
    A_{\mathrm{GC}}
    \exp\left[
        -\frac{(m-m_{\mathrm{TO}})^2}
        {2\sigma_{\mathrm{GCLF}}^2}
    \right].
\end{equation}
Здесь:
\begin{itemize}
    \item \(A_{\mathrm{GC}}\) --- поверхностная плотность
    шаровиков в максимуме функции;
    \item \(m_{\mathrm{TO}}\) --- apparent turnover magnitude,
    положение максимума GCLF;
    \item \(\sigma_{\mathrm{GCLF}}\) --- ширина GCLF.
\end{itemize}
В текущем пайплайне
\begin{equation}
    \sigma_{\mathrm{GCLF}}=1.2~\mathrm{mag}.
\end{equation}

\subsection{Фоновые галактики}

Для фоновых галактик используется степенной закон:
\begin{equation}
    \Phibg(m)=
    A_{\mathrm{bg}}
    10^{\gamma(m-m_{\mathrm{lim}})}.
\end{equation}
Здесь:
\begin{itemize}
    \item \(A_{\mathrm{bg}}\) --- поверхностная плотность
    фоновых галактик на уровне \(m_{\mathrm{lim}}\);
    \item \(\gamma\) --- логарифмический наклон роста числа
    галактик к слабым величинам.
\end{itemize}

Суммарная модель:
\begin{equation}
    \phi(m)=\Phigc(m)+\Phibg(m).
\end{equation}

\section{Переход от звёздной величины к потоку}

Поток источника в единицах изображения определяется как
\begin{equation}
    f(m)=10^{-0.4(m-ZP_{\mathrm{pix}})},
\end{equation}
где \(ZP_{\mathrm{pix}}\) --- AB-нуль-пункт, учитывающий
единицы изображения MJy/sr и площадь пикселя.

Разница в одну звёздную величину соответствует отношению
потоков
\begin{equation}
    \frac{f(m+1)}{f(m)}
    =10^{-0.4}
    \approx0.398.
\end{equation}

\section{Почему в интеграле находится \(f^2\), а не \(f\)}

Средний свет популяции источников был бы равен
\begin{equation}
    F_{\mathrm{mean}}
    =
    \int\phi(m)f(m)\,dm.
\end{equation}
Однако SBF-анализ измеряет дисперсию, то есть мощность
флуктуаций. Для независимых пуассоновских источников
\begin{equation}
    \operatorname{Var}(F)
    =
    \sum_i f_i^2.
\end{equation}
В непрерывном виде:
\begin{equation}
    P_{\mathrm{sources}}
    =
    \int\phi(m)f(m)^2\,dm.
\end{equation}

Квадрат потока приводит к тому, что основной вклад дают
источники непосредственно под пределом обнаружения.
Например, один источник с \(f=10\) имеет мощность
\(f^2=100\), тогда как сто источников с \(f=0.1\) имеют
суммарную мощность всего
\(100\times0.1^2=1\), хотя их суммарный средний поток
одинаков.

\section{Учёт нормировки на модель галактики}

В текущем пайплайне перед FFT используется нормированный
остаток
\begin{equation}
    R_{\mathrm{norm}}(x,y)=
    \frac{
        I(x,y)-I_{\mathrm{model}}(x,y)
    }{
        \sqrt{I_{\mathrm{model}}(x,y)}
    }.
\end{equation}

После этой операции источник с потоком \(f\) имеет
амплитуду
\begin{equation}
    \frac{f}{\sqrt{I_{\mathrm{model}}}},
\end{equation}
а его мощность:
\begin{equation}
    \frac{f^2}{I_{\mathrm{model}}}.
\end{equation}

Поэтому окончательная формула принимает вид
\begin{equation}
\boxed{
    \Prpow =
    \left\langle
        \frac{1}{I_{\mathrm{model}}}
    \right\rangle
    \int_{m_{\mathrm{lim}}}^{m_{\mathrm{lim}}+12}
    \left[
        \Phigc(m)+\Phibg(m)
    \right]
    f(m)^2\,dm
}
\end{equation}

Расшифровка всех обозначений приведена в
таблице~\ref{tab:symbols}.

\begin{table}[ht]
\centering
\caption{Обозначения в формуле \(\Prpow\).}
\label{tab:symbols}
\begin{tabular}{p{0.19\textwidth}p{0.70\textwidth}}
\toprule
Обозначение & Смысл \\
\midrule
\(\Prpow\) & Мощность от незамаскированных компактных источников \\
\(I_{\mathrm{model}}\) & Гладкая модель поверхностной яркости галактики \\
\(\langle I_{\mathrm{model}}^{-1}\rangle\) & Среднее обратной яркости по рабочим пикселям кольца \\
\(m\) & Наблюдаемая AB-звёздная величина источника \\
\(m_{\mathrm{lim}}\) & Предел обнаружения и маскирования каталога \\
\(\Phigc(m)\) & Число шаровых скоплений на пиксель и единицу величины \\
\(\Phibg(m)\) & Число фоновых галактик на пиксель и единицу величины \\
\(f(m)\) & Поток источника в единицах изображения \\
\(f(m)^2\) & Вклад источника в дисперсию, а не в средний свет \\
\bottomrule
\end{tabular}
\end{table}

Верхняя граница \(m_{\mathrm{lim}}+12\) заменяет
бесконечность. Более слабые источники дают пренебрежимо
малый вклад, поскольку
\begin{equation}
    f(m)^2\propto10^{-0.8m}.
\end{equation}

\section{Получение флуктуационной звёздной величины}

После спектрального фита известна величина \(P_0\), а после
анализа каталога --- \(\Prpow\). Тогда
\begin{equation}
    \Pfluc=P_0-\Prpow.
\end{equation}
Флуктуационная звёздная величина:
\begin{equation}
    \overline m =
    -2.5\log_{10}(\Pfluc)
    +ZP_{\mathrm{pix}}.
\end{equation}

Поскольку остатки уже нормированы на
\(\sqrt{I_{\mathrm{model}}}\), \(P_0\) повторно не делится
на среднюю яркость галактики.

Если \(\Prpow\) не вычесть, мощность SBF будет завышена,
\(\overline m\) получится слишком яркой, а расстояние будет
занижено.

Поправка относительно результата без \(\Prpow\):
\begin{equation}
    \Delta\overline m
    =
    -2.5\log_{10}
    \left(
        1-\frac{\Prpow}{P_0}
    \right).
\end{equation}
При \(\Prpow/P_0\ll1\):
\begin{equation}
    \Delta\overline m
    \approx
    1.086\frac{\Prpow}{P_0}.
\end{equation}

Например, при
\begin{equation}
    P_0=1.000,\qquad \Prpow=0.005
\end{equation}
получается
\begin{equation}
    \Pfluc=0.995,\qquad
    \Delta\overline m\approx0.0054~\mathrm{mag}.
\end{equation}

\section{Отличие \(\Prpow\) от \(P_1\)}

\begin{center}
\begin{tabular}{p{0.22\textwidth}p{0.33\textwidth}p{0.33\textwidth}}
\toprule
 & \(\Prpow\) & \(P_1\) \\
\midrule
Физический источник &
Слабые реальные шаровики и фоновые галактики &
Шум изображения \\
Форма по \(k\) &
PSF-подобная, \(E(k)\) &
Плоская в классической модели \\
Где находится &
Внутри измеренного \(P_0\) &
Отдельный член спектрального фита \\
Как определяется &
Каталог и функции светимости &
Одновременно с \(P_0\) из \(P(k)\) \\
\bottomrule
\end{tabular}
\end{center}

Полная идеализированная модель:
\begin{equation}
    P(k)=
    \left(\Psbf+\Prpow\right)E(k)+P_1.
\end{equation}

\section{Алгоритм текущей реализации}

Для каждого рабочего кольца выполняются следующие шаги:
\begin{enumerate}
    \item Из единого каталога выбираются источники, попавшие
    в геометрическую область кольца.
    \item Источники фильтруются по конечной AB-величине и
    положительному потоку.
    \item На интервале от яркого края каталога до
    \(m_{\mathrm{lim}}\) строится гистограмма.
    \item Число источников делится на площадь кольца в пикселях
    и ширину бина по величине.
    \item К гистограмме подгоняется сумма гауссовской GCLF
    и степенной функции фоновых галактик.
    \item Если локальный фит невозможен, используется
    глобальная функция светимости по общей научной области.
    \item Обе функции интегрируются от \(m_{\mathrm{lim}}\)
    до \(m_{\mathrm{lim}}+12\) с весом \(f(m)^2\).
    \item Интеграл умножается на
    \(\langle I_{\mathrm{model}}^{-1}\rangle\).
    \item Получаются \(P_{\mathrm{r,GC}}\),
    \(P_{\mathrm{r,bg}}\) и их сумма \(\Prpow\).
    \item Проверяется условие \(P_0-\Prpow>0\).
    \item Вычисляются \(\Pfluc\) и
    \(\overline m\).
\end{enumerate}

В итоговой таблице каждого кольца сохраняются:
\begin{itemize}
    \item \texttt{P0};
    \item \texttt{Pr\_gc};
    \item \texttt{Pr\_background};
    \item \texttt{Pr};
    \item \texttt{Pr\_over\_P0};
    \item \texttt{P\_fluc};
    \item \texttt{m\_lim};
    \item \texttt{gclf\_turnover\_mag};
    \item \texttt{gclf\_sigma\_mag};
    \item \texttt{lf\_gamma};
    \item \texttt{pr\_note}.
\end{itemize}

\section{Численный масштаб для GO-3055}

В текущем прогоне GO-3055 медианные отношения равны
\begin{align}
    \operatorname{median}
    \left(\frac{\Prpow}{P_0}\right)_{\mathrm{inner}}
    &\approx3.6\times10^{-5},\\
    \operatorname{median}
    \left(\frac{\Prpow}{P_0}\right)_{\mathrm{outer}}
    &\approx1.9\times10^{-3}.
\end{align}
Максимальное значение составляет приблизительно
\begin{equation}
    \max\left(\frac{\Prpow}{P_0}\right)
    \approx6.8\times10^{-3}.
\end{equation}

Следовательно, типичная поправка:
\begin{itemize}
    \item во внутреннем кольце --- существенно меньше
    \(0.001\) mag;
    \item во внешнем кольце --- около \(0.002\) mag;
    \item в максимальном случае --- около \(0.007\) mag.
\end{itemize}

Для GO-3055 поправка мала и не может объяснить изменение
\(\overline m_{150}\) примерно на \(0.3\) mag между старой и
новой версиями пайплайна. Этот сдвиг связан прежде всего с
нормировкой флуктуационного поля, построением \(E(k)\), PSF и
маской.

На расстоянии Coma \(\Prpow\) может стать существенно важнее:
больше шаровых скоплений и фоновых галактик останутся слабее
предела индивидуального обнаружения.

\section{Ограничения и необходимые проверки}

Формула \(\Prpow\) физически обоснована, но результат зависит
от качества входной функции светимости. Основные ограничения
текущей реализации:
\begin{enumerate}
    \item \(m_{\mathrm{lim}}\) оценивается приближённо, без
    полноценного теста полноты искусственными источниками.
    \item Фиксированная ширина GCLF \(1.2\) mag может быть
    неточной для отдельных массивных галактик.
    \item Локальный фит функции светимости может быть
    неустойчив при малом числе источников.
    \item Глобальный fallback уменьшает шум, но может скрыть
    реальное радиальное изменение системы шаровых скоплений.
    \item Пыль, диффузные структуры и артефакты не описываются
    функцией \(\Prpow\).
    \item Слабые foreground-звёзды не выделены отдельной
    компонентой.
\end{enumerate}

Для статьи по GO-3055 достаточно показать малость
\(\Prpow/P_0\) и проверить чувствительность результата к
разумному изменению \(m_{\mathrm{lim}}\). Для Coma потребуется
полноценный тест искусственными источниками.

\section{Формулировка для текста статьи}

\begin{quote}
Измеренная амплитуда PSF-подобной компоненты спектра мощности
\(P_0\) включает как флуктуации неразрешённого звёздного
населения галактики, так и дисперсию потока от шаровых
скоплений и фоновых галактик, оставшихся слабее предела
маскирования. Вклад последних, \(\Prpow\), оценивался
интегрированием суммы гауссовской функции светимости шаровых
скоплений и степенной функции числа фоновых галактик с весом,
равным квадрату потока источника. После приведения к
нормировке остаточного изображения использовалась
флуктуационная мощность
\(\Pfluc=P_0-\Prpow\).
\end{quote}

\section{Источники и связь с проектом}

Формализм соответствует стандартной процедуре коррекции SBF
за неразрешённые внешние источники, описанной, в частности,
в работе Jensen et al. (2015, ApJ, 808, 91). Конкретная
реализация проекта находится в блоке
\texttt{prepare Jensen GCLF + background-galaxy variance model}
ноутбука \texttt{code/sbf-2.ipynb}.

\clearpage
\part{Пересборка пайплайна относительно курсовой работы}

\section{Что именно сравнивается}

В этой части под \textbf{версией курсовой работы} понимается
набор результатов, сохранённый в
\path{code/sbf2_batch_outputs/}, прежде всего в файлах
\path{coursework_calibration_input.csv},
\path{coursework_fit_summary.csv} и
\path{coursework_leave_one_out_residuals.csv}.

Под \textbf{новой Jensen-подобной версией} понимается полный
прогон 14 галактик шаблоном \path{code/sbf-2.ipynb} с
SHA-256
\begin{center}
\texttt{515ad22816f7f621635483938741efc95a48424ade537aac00a5f485a8b6a83a},
\end{center}
сохранённый в
\path{runs/sbf2_go3055/}. Для чисел ниже использованы таблицы
\path{go3055_master_measurements.csv},
\path{go3055_calibration_table.csv},
\path{go3055_fit_summary.csv} и
\path{go3055_leave_one_out_distances.csv} из поддиректории
\path{analysis/tables/}.

Это сравнение двух \emph{реально сохранённых} наборов
результатов. Оно не является сравнением двух графиков,
считанных на глаз. В обоих случаях использованы те же 14
галактик GO-3055, те же входные фильтры F150W и F090W и те
же литературные TRGB-модули расстояния. Изменились способы
обработки изображений, извлечения SBF, измерения цвета и
формирования калибровочной выборки.

\section{Что осталось неизменным}

\begin{itemize}
    \item SBF измеряется по изображениям JWST/NIRCam F150W.
    \item Цветовой показатель строится из F090W и F150W.
    \item Используются готовые продукты уровня \texttt{i2d}.
    \item Основные области измерения --- круговые кольца
    \(8.2''\!-\!16.4''\) и \(16.4''\!-\!32.8''\).
    \item Спектр по-прежнему описывается классической формой
    \(P(k)=P_0E(k)+P_1\).
    \item Вклад неразрешённых источников вычитается как
    \(P_{\mathrm{fluc}}=P_0-P_{\mathrm r}\).
    \item Сохранено симметричное ограничение остаточного
    изображения на уровне \(3.5\sigma\). Это осознанное
    отличие от буквального алгоритма Jensen.
\end{itemize}

\section{Все принципиальные изменения метода}

\footnotesize
\begin{longtable}{
    >{\raggedright\arraybackslash}p{0.14\textwidth}
    >{\raggedright\arraybackslash}p{0.25\textwidth}
    >{\raggedright\arraybackslash}p{0.31\textwidth}
    >{\raggedright\arraybackslash}p{0.20\textwidth}}
\caption{Изменения пайплайна относительно версии курсовой
работы. Последний столбец фиксирует не намерение, а
фактическое следствие для анализа.}
\label{tab:pipeline_changes}\\
\toprule
Этап & Версия курсовой работы & Новая Jensen-подобная версия
& Следствие \\
\midrule
\endfirsthead
\toprule
Этап & Версия курсовой работы & Новая Jensen-подобная версия
& Следствие \\
\midrule
\endhead
\midrule
\multicolumn{4}{r}{Продолжение на следующей странице}\\
\endfoot
\bottomrule
\endlastfoot

Фон &
После начальной оценки применялась пространственная
\texttt{Background2D}-модель. &
Вычитается один скалярный уровень неба. Он оценивается по
углам кадра, профилю \(r^{1/4}\) и остатку после грубой
модели; расхождение оценок сохраняется как систематика. &
Меньше риск вычесть вместе с фоном внешние области
галактики. Полный вклад ошибки фона в новой итоговой таблице
пока не заполнен. \\

Первичная маска &
Источники первоначально искались непосредственно на кадре с
глобальным порогом; дополнительный поиск выполнялся позднее
на остатках. &
Сначала строится предварительная гладкая модель, затем
источники ищутся на остатках с эмпирической дисперсией шума
как функцией локальной яркости модели. &
Собственные SBF реже принимаются за шаровые скопления, а
центр галактики реже попадает в маску как один источник. \\

Каталог и предел маскирования &
Порог детектирования, окончательная маска и граница интеграла
\(P_{\mathrm r}\) не были сведены к одному фотометрическому
пределу. &
Строится единый каталог. Один предел \(m_{\mathrm{cut}}\)
одновременно определяет явно маскируемые источники и нижнюю
границу интеграла \(P_{\mathrm r}\). &
Маска и статистическая поправка теперь описывают
взаимоисключающие части одного каталога. \\

Модель \(P_{\mathrm r}\) &
Использовалась одна объединённая степенная функция
светимости компактных источников. &
Раздельно подгоняются гауссовская GCLF с
\(\sigma_{\mathrm{GCLF}}=1.2\) mag и степенная функция
фоновых галактик. &
Физическая модель приведена к схеме Jensen. Численно
\(P_{\mathrm r}\) остаётся малым и не объясняет общий сдвиг
SBF-величин. \\

Щели и NaN &
Sérsic-модель могла использоваться как рабочее заполнение
больших пропусков и поддержка продолжения модели. &
Заполнение используется только для численной устойчивости
фита изофот. Синтетические пиксели исключаются из FFT и
цветовой фотометрии. &
Щели между детекторами больше не дают искусственный SBF-
сигнал. \\

Full-frame модель &
Модель могла продолжаться за последнюю надёжно измеренную
изофоту. &
На полный кадр переносится только семейство реально
сошедшихся изофот; яркость и геометрия наружу не
экстраполируются. Сохраняется доля покрытия каждого кольца. &
Появились честные предупреждения о неполном покрытии вместо
скрытой синтетической модели. \\

Нормировка остатка &
FFT считался от ненормированного остатка, а \(P_0\) затем
делился на одно среднее значение модели в кольце. &
До FFT используется поле
\((I-I_{\mathrm{model}})/\sqrt{I_{\mathrm{model}}}\). &
Профиль поверхностной яркости входит в измерение до
усреднения. Это главное изменение нормировки и один из
основных кандидатов на сдвиг \(\overline m_{150}\). \\

Спектр \(E(k)\) &
Ожидаемый спектр строился из одного PSF и бинарной маски;
радиальные значения мощности оценивались преимущественно
медианой. &
Для каждого PSF отдельно моделируется поле, свёрнутое с PSF
и умноженное на фактическое окно. Радиальные спектры
усредняются, а ошибки среднего входят в веса фита. &
Форма \(E(k)\) и вклад отдельных частот теперь ближе к
классическому анализу Jensen. \\

Диапазон \(k\) &
Использовались несколько нормированных диапазонов без
однозначной связи с номером Fourier-моды. &
Моды ниже \(k=10\) исключаются явно; основной диапазон
зафиксирован как \(0.04\leq k\leq0.25\), а
\(k_{\min}=0.01\) и \(0.03\) оставлены как диагностика. &
Крупномасштабные ошибки модели меньше влияют на основной
фит; разброс диапазонов остаётся проверкой устойчивости. \\

Подгонка \(P(k)\) &
Использовался невзвешенный линейный МНК. &
Подгонка \(P_0E(k)+P_1\) взвешивается по ошибкам радиальных
средних и повторяется для каждого PSF библиотеки. &
Формальные ошибки лучше отражают неодинаковое число Fourier-
мод в разных интервалах. \\

PSF &
Использовался один синтетический STPSF; OPD выбирался по
времени изменения локального файла. &
OPD выбирается по близости к дате наблюдения. Используются
ближайший WSS OPD и четыре смещения положения по детектору;
при наличии можно добавить эмпирический PSF. &
Для каждого объекта получено пять PSF-реализаций и оценка
PSF-систематики. \\

Цвет &
Кадры приводились к общему размеру, а основной цвет
оценивался через отношение медианных потоков без явной
WCS-перепроекции и согласования PSF. &
F090W перепроецируется на сетку F150W, оба изображения
согласуются по PSF, затем берётся отношение интегральных
потоков в тех же кольцах и по общей валидной маске. &
Цвета стали краснее в среднем на \(0.209\) mag. Старый и
новый цветовые диапазоны вообще не пересекаются. \\

Объединение колец &
Применялось обратное дисперсионное взвешивание; различие
колец входило в ошибку. &
Схема сохранена. Дополнительно записываются покрытие,
используемая доля пикселей и значимость расхождения колец. &
Пять объектов выделены как структурно проблемные, а не
молча смешаны с остальными. \\

Ошибки &
Итоговый бюджет содержал фит \(P(k)\), PSF, \(P_{\mathrm r}\),
фон, поглощение и разброс колец, частично объединённые
эвристически. &
Новая таблица надёжно содержит измерительный, PSF- и
\(P_{\mathrm r}\)-вклады. Поля фоновой и extinction-
систематики пока пусты. &
Новые \(\sigma_{\overline m}\) являются внутренними
ошибками, а не полным публикационным бюджетом. \\

Пакетный запуск &
Результаты разных запусков могли перезаписывать друг друга,
а состояние определялось по файлам и SHA. &
Создан отдельный \path{run_sbf_2_batch.py}: уникальные
директории запусков, табличные статусы, журналы, manifests
артефактов и полностью автономная обработка без сети. &
Для каждого числа можно восстановить шаблон, входные файлы,
PSF, время и выходные продукты. \\

Калибровочный анализ &
Основным был эффективный ``чистый'' набор из 12 объектов,
одна линейная калибровка и LOO-проверка. &
Показываются все 14 объектов и отдельная структурно чистая
подвыборка из 9; Virgo, Fornax и прочие группы разделены,
добавлен тест ступеньки и окраска по светимости. &
Выбросы и кластерная структура больше не скрываются одним
фитом. Прямая цифра разброса стала хуже. \\
\end{longtable}
\normalsize

\section{Новые измерения по всем 14 галактикам}

Обозначим
\begin{equation}
    C = F090W-F150W,
\end{equation}
а разности ``новое минус курсовая'' --- через
\(\Delta\overline m\), \(\Delta C\) и
\(\Delta\overline M\). Поскольку литературный TRGB-модуль
для каждой галактики не менялся,
\begin{equation}
    \Delta\overline M_{150} =
    \Delta\overline m_{150}.
\end{equation}

\begin{landscape}
\scriptsize
\begin{longtable}{llrrrrrr}
\caption{Прямое сравнение наблюдаемой SBF-величины и цвета
для всех галактик. Все величины даны в mag.}
\label{tab:old_new_measurements}\\
\toprule
Галактика & Среда &
\(\overline m_{\rm old}\) &
\(\overline m_{\rm new}\) &
\(\Delta\overline m\) &
\(C_{\rm old}\) &
\(C_{\rm new}\) &
\(\Delta C\) \\
\midrule
\endfirsthead
\toprule
Галактика & Среда &
\(\overline m_{\rm old}\) &
\(\overline m_{\rm new}\) &
\(\Delta\overline m\) &
\(C_{\rm old}\) &
\(C_{\rm new}\) &
\(\Delta C\) \\
\midrule
\endhead
\bottomrule
\endfoot
NGC 1380 & Fornax & 27.887 & 28.120 & \(+0.233\) & 0.413 & 0.522 & \(+0.109\) \\
NGC 1399 & Fornax & 28.125 & 28.186 & \(+0.061\) & 0.438 & 0.578 & \(+0.140\) \\
NGC 1404 & Fornax & 28.004 & 28.212 & \(+0.208\) & 0.477 & 0.567 & \(+0.090\) \\
NGC 1549 & Other  & 27.670 & 27.989 & \(+0.319\) & 0.344 & 0.541 & \(+0.198\) \\
NGC 3379 & Other  & 26.654 & 27.061 & \(+0.407\) & 0.353 & 0.574 & \(+0.221\) \\
NGC 4374 & Virgo  & 27.698 & 27.985 & \(+0.287\) & 0.421 & 0.594 & \(+0.173\) \\
NGC 4406 & Virgo  & 27.487 & 27.878 & \(+0.391\) & 0.306 & 0.637 & \(+0.331\) \\
NGC 4472 & Virgo  & 27.513 & 27.775 & \(+0.262\) & 0.405 & 0.640 & \(+0.235\) \\
NGC 4486 & Virgo  & 27.600 & 27.892 & \(+0.293\) & 0.463 & 0.668 & \(+0.205\) \\
NGC 4552 & Virgo  & 27.618 & 27.897 & \(+0.279\) & 0.458 & 0.598 & \(+0.140\) \\
NGC 4621 & Virgo  & 27.350 & 27.732 & \(+0.382\) & 0.313 & 0.588 & \(+0.275\) \\
NGC 4636 & Virgo  & 27.548 & 27.975 & \(+0.427\) & 0.311 & 0.624 & \(+0.314\) \\
NGC 4649 & Virgo  & 27.569 & 27.975 & \(+0.405\) & 0.326 & 0.641 & \(+0.315\) \\
NGC 4697 & Virgo  & 26.741 & 27.127 & \(+0.386\) & 0.358 & 0.533 & \(+0.175\) \\
\midrule
Среднее & & & & \(+0.310\) & & & \(+0.209\) \\
Медиана & & & & \(+0.306\) & & & \(+0.201\) \\
Минимум--максимум & & & & \(+0.061\ldots+0.427\) &
& & \(+0.090\ldots+0.331\) \\
\end{longtable}
\end{landscape}

Сдвиг \(\overline m_{150}\) имеет один знак у всех 14
объектов. Поэтому это не случайное поведение одной галактики,
а систематическое следствие изменения метода. Главные
кандидаты --- нормировка остатка на
\(\sqrt{I_{\mathrm{model}}}\), новое построение \(E(k)\),
изменившаяся маска и библиотека PSF. Поправка
\(P_{\mathrm r}\) слишком мала: её типичный эффект составляет
тысячные доли звёздной величины.

Старые цвета занимают диапазон
\begin{equation}
    0.306 \leq C_{\mathrm{old}} \leq 0.477,
\end{equation}
а новые ---
\begin{equation}
    0.522 \leq C_{\mathrm{new}} \leq 0.668.
\end{equation}
Диапазоны не пересекаются. Это означает, что старый и новый
наклоны цветовой калибровки нельзя сравнивать как две оценки
одной регрессии при неизменном предикторе. Сначала необходимо
подтвердить нуль-пункт новой согласованной фотометрии
F090W--F150W на нескольких апертурах вручную.

\section{Абсолютные величины и расстояния}

В таблице~\ref{tab:old_new_distances} приведены абсолютные
SBF-величины и остатки leave-one-out:
\begin{equation}
    \Delta\mu_{\mathrm{LOO}}
    = \mu_{\mathrm{SBF,LOO}}-\mu_{\mathrm{TRGB}}.
\end{equation}
Знак ``---'' означает, что галактика не входила в старую
12-объектную LOO-выборку. Последний столбец показывает
разность внутреннего и внешнего колец в новом прогоне.

\begin{landscape}
\scriptsize
\begin{longtable}{lrrrrrrc}
\caption{Абсолютные SBF-величины и LOO-остатки. Все численные
величины, кроме признака чистоты, даны в mag.}
\label{tab:old_new_distances}\\
\toprule
Галактика &
\(\overline M_{\rm old}\) &
\(\overline M_{\rm new}\) &
\(\Delta\overline M\) &
\(\Delta\mu_{\rm LOO,old}\) &
\(\Delta\mu_{\rm LOO,new}\) &
\(|\overline m_{\rm in}-\overline m_{\rm out}|\) &
Структурно чистая \\
\midrule
\endfirsthead
\toprule
Галактика &
\(\overline M_{\rm old}\) &
\(\overline M_{\rm new}\) &
\(\Delta\overline M\) &
\(\Delta\mu_{\rm LOO,old}\) &
\(\Delta\mu_{\rm LOO,new}\) &
\(|\overline m_{\rm in}-\overline m_{\rm out}|\) &
Структурно чистая \\
\midrule
\endhead
\bottomrule
\endfoot
NGC 1380 & -3.521 & -3.288 & \(+0.233\) & \(-0.090\pm0.067\) & \(-0.100\pm0.076\) & 0.045 & нет \\
NGC 1399 & -3.373 & -3.312 & \(+0.061\) & \(+0.038\pm0.112\) & \(-0.145\pm0.062\) & 0.019 & да \\
NGC 1404 & -3.362 & -3.154 & \(+0.208\) & \(+0.004\pm0.067\) & \(+0.036\pm0.076\) & 0.026 & да \\
NGC 1549 & -3.561 & -3.242 & \(+0.319\) & \(-0.051\pm0.059\) & \(-0.047\pm0.078\) & 0.039 & да \\
NGC 3379 & -3.424 & -3.017 & \(+0.407\) & \(+0.096\pm0.063\) & \(+0.173\pm0.109\) & 0.169 & нет \\
NGC 4374 & -3.414 & -3.127 & \(+0.287\) & \(+0.019\pm0.060\) & \(+0.044\pm0.077\) & 0.043 & да \\
NGC 4406 & -3.609 & -3.218 & \(+0.391\) & \(-0.074\pm0.058\) & \(-0.095\pm0.078\) & 0.062 & да \\
NGC 4472 & -3.503 & -3.241 & \(+0.262\) & \(-0.061\pm0.066\) & \(-0.128\pm0.071\) & 0.039 & да \\
NGC 4486 & -3.446 & -3.154 & \(+0.293\) & \(-0.079\pm0.071\) & \(-0.057\pm0.080\) & 0.058 & да \\
NGC 4552 & -3.305 & -3.026 & \(+0.279\) & \(+0.104\pm0.065\) & \(+0.149\pm0.078\) & 0.086 & нет \\
NGC 4621 & -3.471 & -3.089 & \(+0.382\) & --- & \(+0.089\pm0.079\) & 0.071 & да \\
NGC 4636 & -3.522 & -3.095 & \(+0.427\) & \(+0.042\pm0.071\) & \(+0.058\pm0.076\) & 0.042 & да \\
NGC 4649 & -3.440 & -3.034 & \(+0.405\) & \(+0.116\pm0.067\) & \(+0.113\pm0.092\) & 0.120 & нет \\
NGC 4697 & -3.583 & -3.197 & \(+0.386\) & --- & \(+0.016\pm0.095\) & 0.120 & нет \\
\end{longtable}
\end{landscape}

Структурно проблемными в новой диагностике обозначены
NGC~1380, NGC~3379, NGC~4552, NGC~4649 и NGC~4697.
Причины включают увеличенное расхождение колец, неполное
покрытие внешнего кольца моделью или уменьшенную долю
реально используемых пикселей. Это не означает, что все пять
галактик физически непригодны для SBF; это означает, что их
нельзя использовать как безусловно чистые калибраторы без
просмотра промежуточных FITS.

\subsection{Итоговые расстояния в Мпк}

В таблице~\ref{tab:go3055_distances_mpc} собраны расстояния
TRGB, leave-one-out оценки SBF для четырёх проверяемых
цветовых законов и внешнее сравнение с реконструированными
HST/F160W SBF-расстояниями Jensen et al. (2015). Значения
Jensen существуют только для пяти пересекающихся галактик;
для остальных объектов в исходной таблице стоит
\texttt{NaN}. Ошибки SBF в этой версии являются внутренними:
они ещё не включают неопределённость коэффициентов
калибровки, неба и галактического поглощения.

\begin{landscape}
\begin{table}[p]
\centering
\caption{Расстояния до галактик GO-3055 в Мпк. Для SBF
показаны leave-one-out результаты постоянной, линейной,
квадратичной и ломаной цветовых моделей.}
\label{tab:go3055_distances_mpc}
\begingroup
\catcode`\_=12
\resizebox{\linewidth}{!}{%
\begin{tabular}{llrrrrrrrrrrrr}
\toprule
galaxy & environment & D_trgb_mpc & sigma_D_trgb_mpc & D_sbf_linear_mpc & sigma_D_sbf_linear_mpc & D_sbf_quadratic_mpc & sigma_D_sbf_quadratic_mpc & D_sbf_broken_mpc & sigma_D_sbf_broken_mpc & D_sbf_constant_mpc & sigma_D_sbf_constant_mpc & D_jensen_f160w_mpc & sigma_D_jensen_f160w_mpc \\
\midrule
NGC 1380 & Fornax & 19.125 & 0.220 & 18.052 & 0.794 & 18.462 & 0.658 & 18.531 & 0.652 & 17.972 & 0.692 & 20.120 & 1.191 \\
NGC 1399 & Fornax & 19.934 & 0.239 & 18.564 & 0.693 & 18.178 & 0.515 & 18.002 & 0.468 & 18.544 & 0.710 & 20.298 & 1.185 \\
NGC 1404 & Fornax & 18.759 & 0.216 & 18.950 & 0.788 & 18.722 & 0.677 & 18.679 & 0.666 & 18.876 & 0.797 & 19.841 & 1.223 \\
NGC 1549 & Other & 17.628 & 0.211 & 17.133 & 0.751 & 17.180 & 0.609 & 17.278 & 0.608 & 16.997 & 0.705 & NaN & NaN \\
NGC 3379 & Other & 10.366 & 0.134 & 11.148 & 0.408 & 11.046 & 0.353 & 11.036 & 0.349 & 11.090 & 0.421 & NaN & NaN \\
NGC 4374 & Virgo & 16.688 & 0.207 & 16.933 & 0.703 & 16.697 & 0.609 & 16.596 & 0.596 & 16.894 & 0.719 & NaN & NaN \\
NGC 4406 & Virgo & 16.565 & 0.191 & 15.851 & 0.659 & 15.832 & 0.546 & 15.943 & 0.551 & 16.073 & 0.671 & NaN & NaN \\
NGC 4472 & Virgo & 15.966 & 0.199 & 15.051 & 0.632 & 15.140 & 0.527 & 15.206 & 0.527 & 15.342 & 0.639 & 16.190 & 0.983 \\
NGC 4486 & Virgo & 16.188 & 0.186 & 15.842 & 0.846 & 16.909 & 0.695 & 16.399 & 0.675 & 16.223 & 0.765 & NaN & NaN \\
NGC 4552 & Virgo & 15.297 & 0.197 & 16.295 & 0.608 & 16.111 & 0.542 & 16.054 & 0.538 & 16.274 & 0.632 & NaN & NaN \\
NGC 4621 & Virgo & 14.595 & 0.175 & 15.138 & 0.607 & 14.948 & 0.532 & 14.873 & 0.525 & 15.101 & 0.623 & NaN & NaN \\
NGC 4636 & Virgo & 16.368 & 0.196 & 16.479 & 0.709 & 16.343 & 0.598 & 16.466 & 0.594 & 16.596 & 0.708 & NaN & NaN \\
NGC 4649 & Virgo & 15.915 & 0.183 & 16.826 & 0.720 & 16.927 & 0.552 & 17.052 & 0.532 & 16.923 & 0.663 & 15.856 & 0.984 \\
NGC 4697 & Virgo & 11.609 & 0.150 & 11.527 & 0.562 & 11.922 & 0.427 & 11.962 & 0.422 & 11.365 & 0.482 & NaN & NaN \\
\bottomrule
\end{tabular}
%
}
\endgroup
\end{table}
\end{landscape}

\section{Сравнение калибровок}

Старая основная калибровка строилась по 12 объектам:
\begin{equation}
    \overline M_{150}
    = -3.454
    +1.165\left(C-0.400\right).
\end{equation}
Новый фит по всем 14 объектам:
\begin{equation}
    \overline M_{150}
    = -3.170
    +0.714\left(C-0.591\right).
\end{equation}
Новый фит по девяти структурно чистым объектам:
\begin{equation}
    \overline M_{150}
    = -3.188
    +0.364\left(C-0.591\right).
\end{equation}

\begin{table}[htbp]
\centering
\small
\caption{Сводка старой и новых линейных калибровок
\(\overline M_{150}=A+B(C-C_0)\). Новый LOO-WRMS рассчитан
непосредственно из сохранённой таблицы расстояний.}
\label{tab:fit_comparison}
\begin{tabular}{lrrrrrr}
\toprule
Вариант & \(N\) & \(C_0\) & \(A\) & \(B\) &
\(\sigma_{\mathrm{int}}\) & LOO-WRMS \\
\midrule
Курсовая, принятая & 12 & 0.400 & -3.454 & 1.165 & 0.049 & 0.073 \\
Новая, все объекты & 14 & 0.591 & -3.170 & 0.714 & 0.072 & 0.100 \\
Новая, структурно чистая & 9 & 0.591 & -3.188 & 0.364 & 0.062 & 0.084 \\
\bottomrule
\end{tabular}
\end{table}

Для старой калибровки формальные ошибки коэффициентов были
\(0.013\) и \(0.209\). Для нового полного фита 16-й и 84-й
процентили bootstrap равны
\([-3.189,-3.143]\) для \(A\) и
\([0.237,1.162]\) для \(B\). Для структурно чистого фита они
равны \([-3.214,-3.158]\) и
\([-0.193,0.850]\), соответственно. Следовательно, наклон
структурно чистой выборки статистически совместим с нулём.

Новая реализация \emph{не улучшила} наблюдаемый разброс:
LOO-WRMS вырос с \(0.073\) до \(0.100\) mag для полного
набора. Даже после исключения пяти структурно проблемных
объектов он равен \(0.084\) mag. Это не доказывает, что новая
физика хуже. Скорее, старая точность была частично
оптимистичной, а новая обработка раскрыла зависимость от
маски, модели, цвета и структуры колец.

\section{Проверка возможной ступеньки Virgo--Fornax}

Для 12 галактик, отнесённых к Virgo или Fornax, был
подогнан дополнительный сдвиг Fornax:
\begin{equation}
    \overline M_{150}
    = A+B(C-C_0)+\Delta_{\mathrm{Fornax}}.
\end{equation}
Получены
\begin{align}
    \Delta_{\mathrm{Fornax,MLE}} &= -0.125~\mathrm{mag},\\
    \operatorname{median}
    \left(\Delta_{\mathrm{Fornax,boot}}\right)
    &= -0.140~\mathrm{mag},\\
    P_{16}\ldots P_{84}
    &= -0.209\ldots-0.065~\mathrm{mag}.
\end{align}
Однако
\begin{equation}
    \Delta\mathrm{AICc}
    = \mathrm{AICc}_{\mathrm{step}}
    -\mathrm{AICc}_{\mathrm{no\ step}}
    = +0.91.
\end{equation}
То есть дополнительная ступенька \emph{не предпочитается}
по AICc, несмотря на ненулевой bootstrap-интервал. На 12
объектах это следует называть намёком или диагностикой, но
не обнаружением двух независимых нуль-пунктов.

\section{Проблема отрицательного \(P_1\)}

В новом прогоне \(P_1<0\) получен \textbf{для всех 14
галактик и в обоих кольцах}. Алгебраически отрицательный
свободный коэффициент допустим, но физически \(P_1\)
интерпретируется как мощность белого шума и не должен быть
отрицательным.

Массовость эффекта исключает объяснение одной неудачной
галактикой. Наиболее вероятная причина --- несоответствие
плоского шумового члена реальному коррелированному шуму
\texttt{i2d}:
\begin{equation}
    P(k)=P_0E(k)+P_1
\end{equation}
слишком жёстко заставляет всю неплоскую часть спектра
распределяться между \(E(k)\) и отрицательной константой.
До публикации необходимо параллельно проверить модель
\begin{equation}
    P(k)=P_0E(k)+P_1N(k),
\end{equation}
где \(N(k)\) --- измеренная форма спектра шума пустого неба
или эквивалентного контрольного кадра. Классический фит
следует сохранить для прямого сравнения с Jensen.

\section{Что уже можно считать новым результатом}

\begin{enumerate}
    \item Полностью воспроизводимо обработаны все 14
    галактик GO-3055 по единому Jensen-подобному алгоритму.
    \item Показан устойчивый систематический сдвиг
    \(\overline m_{150}\): среднее \(+0.310\) mag, диапазон
    \(+0.061\ldots+0.427\) mag.
    \item Переопределение цветовой фотометрии дало средний
    сдвиг \(+0.209\) mag; старый и новый диапазоны цвета не
    пересекаются.
    \item Новая калибровка по всем объектам имеет
    \(\sigma_{\mathrm{int}}=0.072\) mag и LOO-WRMS
    \(0.100\) mag; структурно чистая подвыборка ---
    \(0.062\) и \(0.084\) mag.
    \item Выделены пять объектов с выраженной структурной
    систематикой колец.
    \item Получен намёк на сдвиг Fornax относительно Virgo
    около \(-0.13\) mag, но модель со ступенькой не
    предпочитается по AICc.
    \item Показано, что \(P_{\mathrm r}\) мало:
    медиана \(P_{\mathrm r}/P_0\) равна
    \(3.6\times10^{-5}\) во внутреннем и
    \(1.9\times10^{-3}\) во внешнем кольце; максимум
    \(6.8\times10^{-3}\).
    \item Обнаружена общая для выборки проблема:
    \(P_1<0\) в 28 из 28 измеренных колец.
\end{enumerate}

\section{Что пока не является итогом}

\begin{enumerate}
    \item Нельзя выбирать между старой и новой калибровкой
    только по меньшему разбросу. Старая версия методически
    дальше от Jensen, а новая пока имеет нерешённую проблему
    шумовой модели.
    \item Нельзя интерпретировать новый цветовой наклон до
    независимой проверки фотометрического нуль-пункта,
    перепроекции и PSF matching хотя бы для NGC~1380,
    NGC~1404 и одной Virgo-галактики.
    \item Нельзя называть внутренние ошибки новой таблицы
    полным бюджетом: вклад фона, поглощения и часть
    систематик ещё не заполнены.
    \item Нельзя объявлять ступеньку Virgo--Fornax
    обнаруженной: \(\Delta\mathrm{AICc}=+0.91\).
    \item Нельзя публиковать классический \(P_0E(k)+P_1\)
    как физически полностью согласованный фит, пока
    отрицательный \(P_1\) остаётся у всех объектов.
\end{enumerate}

\section{Формулировка для научного руководителя}

\begin{quote}
После курсовой работы пайплайн был приведён ближе к процедуре
Jensen: изменены оценка фона, последовательность построения
маски, нормировка SBF-поля, модель \(P_{\mathrm r}\), выбор и
ансамблирование PSF, радиальное усреднение спектра и
фотометрия цвета в согласованных по WCS и PSF кольцах.
Перезапуск тех же 14 галактик дал систематически более слабые
SBF, в среднем на \(0.310\) mag, и более красные цвета, в
среднем на \(0.209\) mag. Новая LOO-точность составляет
\(0.100\) mag для полного набора и \(0.084\) mag для девяти
структурно чистых объектов против \(0.073\) mag в курсовой
версии. Есть намёк на сдвиг Fornax относительно Virgo около
\(-0.13\) mag, но AICc не требует ступенчатой модели.
Критическая нерешённая проблема --- отрицательный белый
шум \(P_1\) во всех 28 кольцах, поэтому следующий
обязательный тест --- параллельный фит с эмпирической формой
шума \(N(k)\).
\end{quote}

\end{document}
